% Part: propositional-logic
% Chapter: propositional-logic
% Section: preliminaries

\documentclass[../../../include/open-logic-section]{subfiles}

\begin{document}

\olfileid{pl}{syn}{pre}

\olsection{पूर्वतयारी}

\begin{thm}[\emph{!!{formula}s वरील विगमनाचे तत्त्व}]
  \ollabel{thm:induction}  
  एखादा गुणधर्म~$P$ सर्व आण्विक !!{formula}s साठी लागू होत असेल आणि
  पुढील अटी पूर्ण करत असेल:
  \begin{enumerate}
    \tagitem{prvNot}{तो~$!A$ साठी लागू असताना $\lnot !A$ साठीही
      लागू होतो;}{ }
    \tagitem{prvAnd}{तो~$!A$ आणि~$!B$ यांच्यासाठी लागू असताना
      $(!A \land !B)$ साठीही लागू होतो;}{ }
    \tagitem{prvOr}{तो~$!A$ आणि~$!B$ यांच्यासाठी लागू असताना
      $(!A \lor !B)$ साठीही लागू होतो;}{ }
    \tagitem{prvIf}{तो~$!A$ आणि~$!B$ यांच्यासाठी लागू असताना
      $(!A \lif !B)$ साठीही लागू होतो;}{ }
    \tagitem{prvIff}{तो~$!A$ आणि~$!B$ यांच्यासाठी लागू असताना
      $(!A \liff !B)$ साठीही लागू होतो;}{ }
  \end{enumerate}
  तर $P$ सर्व !!{formula}s साठी लागू होतो.
\end{thm}

\begin{proof}
  गुणधर्म~$P$ असलेल्या सर्व !!{formula}s चा संग्रह~$S$ असू द्या.
  स्पष्टपणे $S \subseteq \Frm[L_0]$. \olref[fml]{defn:formulas} मधील
  सर्व अटी $S$~पूर्ण करतो: त्यात सर्व आण्विक !!{formula}s आहेत आणि
  तो !!{operator}s खाली बंद आहे.  $\Frm[L_0]$ हा असा सर्वांत लहान
  वर्ग आहे; म्हणून $\Frm[L_0] \subseteq S$. त्यामुळे $\Frm[L_0] = S$,
  आणि प्रत्येक सूत्राला गुणधर्म~$P$ आहे.
\end{proof}

\begin{prop}\ollabel{prop:balanced}
   $\Frm[L_0]$ मधील प्रत्येक !!{formula} \emph{संतुलित} आहे; म्हणजे
   त्यात जेवढे डावे कंस आहेत तेवढेच उजवे कंस आहेत.
\end{prop}

\begin{prob}
\olref[pl][syn][pre]{prop:balanced} सिद्ध करा.
\end{prob}

\begin{prop} \ollabel{prop:noinit}
कोणत्याही !!a{formula} चा कोणताही उचित आरंभीचा खंड !!a{formula} नसतो.
\end{prop}

\begin{prob}
  \olref[pl][syn][pre]{prop:noinit} सिद्ध करा.
\end{prob}

\begin{prop}[एकमेव वाचनीयता]
$ \Frm[L_0]$ मधील कोणत्याही !!{formula}~$!A$ चे पुढीलपैकी नेमक्या एका
रूपात रचनाविश्लेषण करता येते:
\begin{enumerate}
\tagitem{prvFalse}{$\lfalse$.}{ }

\tagitem{prvTrue}{$\ltrue$.}{ }

\item एखाद्या $\Obj p_n \in  \PVar$ साठी $\Obj p_n$.
  
\tagitem{prvNot}{एखाद्या !!{formula}~$!B$ साठी $\lnot !B$.}{ }

\tagitem{prvAnd}{काही !!{formula}s $!B$ आणि~$!C$ यांच्यासाठी
  $(!B \land !C)$.}{ }

\tagitem{prvOr}{काही !!{formula}s $!B$ आणि~$!C$ यांच्यासाठी
  $(!B \lor !C)$.}{ }

\tagitem{prvIf}{काही !!{formula}s $!B$ आणि~$!C$ यांच्यासाठी
  $(!B \lif !C)$.}{ }

\tagitem{prvIff}{काही !!{formula}s $!B$ आणि~$!C$ यांच्यासाठी
  $(!B \liff !C)$.}{ }
\end{enumerate}
शिवाय, हे रचनाविश्लेषण \emph{एकमेव} असते.
\end{prop}

\begin{proof}
$!A$ वरील विगमनाने सिद्ध करू. उदाहरणार्थ, $!A$ चे $(!B \lif !C)$ आणि
$(!B' \lif !C')$ अशी दोन भिन्न रचनाविश्लेषणे आहेत असे समजा. मग $!B$
आणि $!B'$ एकच असले पाहिजेत (अन्यथा त्यांपैकी एक हा दुसऱ्याचा उचित
आरंभीचा खंड असेल); म्हणून $!A$ ची ही दोन रचनाविश्लेषणे भिन्न असतील,
तर त्याचे कारण असेच असले पाहिजे की $!C$ आणि $!C'$ ही एकाच
चिन्हमालेची भिन्न रचनाविश्लेषणे आहेत; पण विगमन गृहीतकानुसार ते अशक्य
आहे.
\end{proof}


\begin{defn}[एकरूप आदेशन]
$!A$ आणि $!B$ ही !!{formula}s असतील आणि $\Obj p_i$ हे !!{propositional
variable} असेल, तर $!A$ मधील $\Obj p_i$ ची प्रत्येक
आस्थिति $!B$ च्या एका आस्थितीने बदलल्यावर मिळणारा परिणाम
$\Subst{!A}{!B}{\Obj p_i}$ ने दर्शवला जातो; त्याचप्रमाणे,
$\Obj p_1$, \dots,~$\Obj p_n$ यांचे !!{formula}s $!B_1$,
\dots,~$!B_n$ यांनी एकाच वेळी केलेले आदेशन
$\SSubst{!A}{\subst{!B_1}{\Obj p_1},\dots,\subst{!B_n}{\Obj p_n}}$
ने दर्शवले जाते.
\end{defn}

\begin{prob} खालील पाच !!{formula}s पैकी प्रत्येकासाठी, ते
 !!{formula} \( \Subst{!A}{!B}{\Obj p_i} \) या आदेशनाच्या रूपात व्यक्त
 करता येते का ते ठरवा; येथे \( !A \) हे अनुक्रमे (i) \( \Obj p_0 \);
 (ii) \( ( \lnot \Obj p_0 \land \Obj p_1) \); आणि (iii)
 \( ( ( \lnot \Obj p_0 \lif \Obj p_1 ) \land \Obj p_2 ) \) आहे.
 प्रत्येक बाबतीत संबद्ध आदेशन स्पष्ट करा.
  \begin{enumerate}
    \item \( \Obj p_1 \) 
    \item \( ( \lnot \Obj p_0 \land \Obj p_0 ) \)
    \item \( ( ( \Obj p_0 \lor \Obj p_1 ) \land \Obj p_2 )  \)
    \item \( \lnot ( ( \Obj p_0 \lif \Obj p_1 ) \land \Obj p_2 ) \)
    \item \( (( \lnot ( \Obj p_0 \lif \Obj p_1 ) \lif ( \Obj p_0 \lor \Obj p_1 )) \land \lnot ( \Obj p_0 \land \Obj p_1 )) \)
  \end{enumerate}
\end{prob}

\begin{prob}
$\Subst{!A}{!B}{p}$ ची विगमनाधारित, गणितीयदृष्ट्या काटेकोर व्याख्या
द्या.
\end{prob}

\end{document}
